\documentclass[aps,prl,twocolumn,showpacs,superscriptaddress,nofootinbib]{revtex4-2}

\usepackage{amsmath,amssymb}
\usepackage{graphicx}
\usepackage{booktabs}
\usepackage{siunitx}
\usepackage{hyperref}

% Canonical counts pipeline (Phase 5v: freshness gate).
% Auto-generated by update_counts.py — 2026-04-24T19:44:03.140400
% Do not edit manually. Run: uv run python scripts/update_counts.py
\newcommand{\totaltheorems}{3993}
\newcommand{\substantivetheorems}{3882}
\newcommand{\placeholdertheorems}{111}
\newcommand{\axiomcount}{1}
\newcommand{\sorrycount}{0}
\newcommand{\sorrypercent}{0.0}
\newcommand{\leanmodules}{167}
\newcommand{\leandefinitions}{3297}
\newcommand{\pythonmodules}{68}
\newcommand{\testfiles}{59}
\newcommand{\totaltests}{2836}
\newcommand{\figurecount}{111}
\newcommand{\notebookcount}{54}
\newcommand{\papercount}{19}
\newcommand{\aristotleproved}{322}
\newcommand{\aristotleruns}{44}
% --- Paper 18 / Phase 5t per-section counts (FermiHubbardDimer.lean) ---
\newcommand{\fhdTotal}{143}
\newcommand{\fhdWaveTwo}{13}
\newcommand{\fhdWaveThree}{6}
\newcommand{\fhdWaveFour}{22}
\newcommand{\fhdWaveFive}{22}
\newcommand{\fhdWaveSixABC}{38}
\newcommand{\fhdWaveSixExtra}{15}
\newcommand{\fhdWaveSeven}{19}
\newcommand{\fhdWaveEight}{8}


\newcommand{\WallIntact}{\mathsf{WallIntact}}
\newcommand{\WallCracked}{\mathsf{WallCracked}}
\newcommand{\EWBGViable}{\mathsf{EWBGViable}}
\newcommand{\Zsixteen}{\mathbb{Z}_{16}}

\begin{document}

\title{Electroweak baryogenesis is doubly forbidden in the Standard
Model: a Lean-formalized two-pillar bridge to the chirality wall}

\author{John G.\ Roehm}
\affiliation{NetRxn Foundation}

\date{\today}

\begin{abstract}
We formalize the verdict that electroweak baryogenesis is forbidden in
the Standard Model under \emph{two independent} obstructions:
\textit{(i)} the chirality-wall obstruction at the gauging step, and
\textit{(ii)} the crossover nature of the SM electroweak phase
transition. The first reflects the Wang $\Zsixteen$ anomaly
classification~\cite{WangWanYou2022}: the SM-no-$\nu_R$ fermion content
gives $3 \times 15 = 45 \equiv 13 \not\equiv 0 \pmod{16}$, blocking
non-perturbative lattice gauging~\cite{Wang2020} and depriving the
sphaleron rate of a first-principles regulator. The second is a
lattice result: the resummed-thermal-corrections endpoint at
$m_H = 72.4 \pm 1.7$~GeV (Csikor, Fodor, and
Heitger~\cite{CsikorFodorHeitger1999}, refining
Kajantie--Laine--Rummukainen--Shaposhnikov~\cite{KLRS1996}) sits well
below the PDG value $m_H = 125.20$~GeV~\cite{PDG2024}; the SM
electroweak transition is a crossover, blocking sphaleron decoupling
(Sakharov's third condition).

We ship \texttt{EWBaryogenesisChiralityWall.lean}: $16$ theorems
formalizing the sphaleron-suppression structural form, the chirality-
wall predicate \WallIntact{} / \WallCracked{} over $\Zsixteen$,
substantive cross-bridges to
\texttt{Z16AnomalyComputation.three\_gen\_anomalous} and
\texttt{sm\_anomaly\_with\_nu\_R}, the compound \EWBGViable{}
predicate, the load-bearing biconditional
$\neg \EWBGViable{} \leftrightarrow (\WallIntact \vee \neg
\mathsf{IsBaryogenesisViable})$, and the punchline
\texttt{sm\_no\_nu\_R\_ewbg\_doubly\_forbidden}: under the
KLRS-crossover hypothesis, the SM-as-is fails \emph{both} conditions
simultaneously. The key insight is that the two failure modes are
\emph{independent}: SM-no-$\nu_R$ fails the wall regardless of the
transition; SM with three right-handed neutrinos cracks the wall but
fails the transition. Either failure suffices to forbid SM EWBG. The
formalization invokes Phase 5z.3
\texttt{crossover\_excludes\_baryogenesis} and Phase 5e
\texttt{three\_gen\_anomalous} in proof bodies; it is zero-sorry,
introduces zero new axioms, and is consistent with the existing
project axiom budget.
\end{abstract}

\pacs{12.15.-y, 11.30.-j, 98.80.Cq, 11.10.Wx}

\maketitle

%% ================================================================
\section{Introduction}
\label{sec:intro}

The Standard Model has two well-known obstructions to electroweak
baryogenesis (EWBG). The first is dynamical: insufficient CP violation
in the Cabibbo--Kobayashi--Maskawa phase, by approximately ten orders
of magnitude relative to what the observed baryon asymmetry
requires~\cite{Gavela1994,Huet1994}. The second is thermodynamic: the
electroweak phase transition is a crossover for the physical Higgs
mass~\cite{KLRS1996,CsikorFodorHeitger1999}, leaving Sakharov's third
condition (departure from thermal equilibrium) unrealized. We formalize
these results and add a third obstruction --- the
\emph{chirality-wall obstruction} from non-perturbative
$\Zsixteen$-anomaly classification~\cite{WangWanYou2022,Wang2020} ---
inside a single Lean module that bridges the chirality-wall pillar
(Phases 5a, 5e, 5p) and the phase-transition pillar (Phase 5z.3).

The two obstructions our paper formalizes act independently. The
SM-no-$\nu_R$ fermion content carries a $\Zsixteen$ anomaly of $13$ mod
$16$ (= $-3$); cancellation is required for the gauging step
(Phase 5p \texttt{gauging\_requires\_z16\_cancellation}), so the
SM-as-is is non-perturbatively obstructed before the question of
transition order even applies. Adding three right-handed neutrinos
brings the anomaly to $48 \equiv 0$ mod $16$, cracking the wall, but
under the KLRS hypothesis the resummed transition is still a crossover
and the exclusion stands.

Our load-bearing theorem is the biconditional
\begin{equation}
\label{eq:biconditional}
\neg \EWBGViable{}(a, p, c)
\;\;\Leftrightarrow\;\;
\bigl( \WallIntact{}(a) \;\vee\;
\neg\, \mathsf{IsBaryogenesisViable}(p, c) \bigr),
\end{equation}
where $a \in \Zsixteen$ is the fermion-content anomaly,
$p \in \mathsf{EWFiniteTParams}$ encodes the finite-$T$ effective
potential parameters, and $c \in \mathbb{R}$ is the Cohen--Kaplan--Nelson
sphaleron-decoupling threshold~\cite{CohenKaplanNelson1993}. The proof
is a de Morgan unfolding with case analysis on $a = 0$ and is
substantive: the LHS-to-RHS direction requires constructive case
splits (mere classical de Morgan would not suffice in the dependent-
type setting). The corollary
\texttt{sm\_no\_nu\_R\_ewbg\_doubly\_forbidden} packages the two
independent failure modes for SM-as-is (intact wall) and
SM-with-$3\nu_R$ (crossover under KLRS) as a conjunction; this is the
formal statement of ``SM EWBG is doubly forbidden.''

The remainder of this paper is organized as follows. Section
\ref{sec:sphaleron} introduces the sphaleron-suppression structural
form and the Cohen--Kaplan--Nelson decoupling predicate. Section
\ref{sec:wall} formalizes the chirality-wall predicate and its
substantive cross-bridges to
\texttt{Z16AnomalyComputation}. Section \ref{sec:bridge} states the
compound \EWBGViable{} and the load-bearing biconditional
\eqref{eq:biconditional}. Section \ref{sec:sm-verdict} applies the
biconditional to the SM with and without right-handed neutrinos under
the KLRS hypothesis and concludes with the doubly-forbidden punchline.
Section \ref{sec:formal-verification} surveys the formalization.
Section \ref{sec:outlook} dispatches baryogenesis to leptogenesis or
BSM mechanisms.

%% ================================================================
\section{Sphaleron Suppression and Decoupling}
\label{sec:sphaleron}

We model the dimensionless Boltzmann suppression of the sphaleron rate
in the broken phase as
\begin{equation}
\label{eq:sphaleron-suppression}
\Sigma(v, T) = e^{-v/T},
\end{equation}
the load-bearing structural form embedded in the
Klinkhamer--Manton sphaleron rate $\Gamma_\mathrm{sphal} \sim
(\alpha_W T)^4 e^{-E_\mathrm{sph}/T}$ with sphaleron energy
$E_\mathrm{sph} \sim 4\pi v / g_W$. The Lean theorem
\texttt{sphaleronSuppression\_pos} establishes
$\Sigma(v, T) > 0$ at any finite $(v, T)$, and
\texttt{sphaleronSuppression\_le\_one} gives the broken-phase bound
$\Sigma \le 1$ when $v/T \ge 0$. Sphaleron decoupling is the
predicate \texttt{SphaleronsDecoupled} encoding $v(T_c)/T_c > c$
with $c = 1$ (Cohen--Kaplan--Nelson~\cite{CohenKaplanNelson1993}).
The threshold is identical to the strong-first-order condition
$E/(\lambda T_c) > c$ used by
\texttt{EWPhaseTransition.IsBaryogenesisViable} (Phase 5z.3).

%% ================================================================
\section{The Chirality-Wall Predicate}
\label{sec:wall}

The chirality wall is a non-perturbative obstruction at the gauging
step: lattice formulations of chiral gauge theory require
$\Zsixteen$-anomaly cancellation~\cite{WangWanYou2022,Wang2020}, an
anomaly-matching constraint formalized as
\texttt{GaugingStep.gauging\_requires\_z16\_cancellation} in Phase 5p.
We define
\begin{align}
\label{eq:wall-intact}
\WallIntact{}(a) &\equiv (a \ne 0 \in \Zsixteen), \\
\label{eq:wall-cracked}
\WallCracked{}(a) &\equiv (a = 0 \in \Zsixteen).
\end{align}
\texttt{wall\_intact\_xor\_cracked} and \texttt{wall\_intact\_or\_cracked}
establish disjointness and exhaustivity over $\Zsixteen$.

The substantive cross-bridges to \texttt{Z16AnomalyComputation}
identify the wall predicate with the project's existing $\Zsixteen$
classification:
\begin{itemize}
\item \texttt{sm\_no\_nu\_R\_wall\_intact} consumes
\texttt{three\_gen\_anomalous} ($45 \ne 0 \in \Zsixteen$) directly
to produce the wall-intact verdict for the SM-no-$\nu_R$ fermion
content (3 generations $\times$ 15 components per generation, per
\texttt{SMFermionData.total\_components\_without\_nu\_R}).
\item \texttt{sm\_no\_nu\_R\_anomaly\_eq\_neg\_three} (substantive
cross-bridge) composes the natural-representative identity
$45 \equiv 13 \pmod{16}$ with the upstream
\texttt{three\_gen\_is\_neg3} (Phase 5e) to identify the SM-no-$\nu_R$
$\Zsixteen$ anomaly with the canonical $-3$ form used in Wang's
classification.
\item \texttt{sm\_single\_gen\_with\_nu\_R\_wall\_cracks} consumes
\texttt{sm\_anomaly\_with\_nu\_R} ($16 \equiv 0 \in \Zsixteen$) to
produce the cracked-wall verdict at the single-generation
right-handed-neutrino-completed fermion content (16 components per
generation).
\item \texttt{sm\_with\_3nu\_R\_wall\_cracks} extends to three
generations: $48 \equiv 0 \pmod{16}$.
\end{itemize}

The substantive content is the \emph{identification} of the
chirality-wall obstruction with the $\Zsixteen$-anomaly nonzero
condition --- not merely the modular arithmetic.

%% ================================================================
\section{Bridge Theorem and Load-Bearing Biconditional}
\label{sec:bridge}

The compound \EWBGViable{} predicate
\begin{equation}
\label{eq:viability}
\EWBGViable{}(a, p, c) \;\equiv\; \WallCracked{}(a) \wedge
\mathsf{IsBaryogenesisViable}(p, c)
\end{equation}
expresses the two-pillar EWBG necessary condition: cracked wall
\emph{and} sphaleron-decoupling-strength first-order transition. The
elementary bridge theorems are
\begin{itemize}
\item \texttt{ewbg\_forbidden\_if\_wall\_intact}: a non-zero
$\Zsixteen$ anomaly forbids EWBG independent of transition order.
\item \texttt{ewbg\_forbidden\_if\_transition\_crossover}: a crossover
transition forbids EWBG independent of wall integrity. Substantive
cross-bridge --- the proof body invokes
\texttt{EWPhaseTransition.crossover\_excludes\_baryogenesis} from
Phase 5z.3.
\end{itemize}

The load-bearing correctness-push is the biconditional
\eqref{eq:biconditional}, formalized as
\texttt{ewbg\_forbidden\_iff\_wall\_intact\_or\_not\_viable}: every
breakdown of EWBG must be attributed to either the wall or the phase
transition, with no third failure mode in this two-pillar
decomposition. The proof is non-trivial: it requires a constructive
case split on $a = 0$ (decidable on $\Zsixteen$). The structural
content makes failure-mode attribution exhaustive in a way that
either elementary bridge alone does not.

%% ================================================================
\section{The SM Verdict: Doubly Forbidden}
\label{sec:sm-verdict}

The KLRS lattice analysis~\cite{KLRS1996,CsikorFodorHeitger1999}
establishes that, after thermal-loop resummation, the SM electroweak
transition is a crossover for $m_H > 72.4 \pm 1.7$~GeV. With the PDG
observed $m_H = 125.20$~GeV~\cite{PDG2024}, the SM is well above the
endpoint. The Lean theorem
\texttt{sm\_klrs\_overshoot\_ratio\_gt\_threshold}
encodes the explicit bound $1.5 < 125.20 / 72.4 \approx 1.73$ via
\texttt{norm\_num}, providing the load-bearing physics input for the
tracked-Prop hypothesis $H_\mathrm{KLRS}$:
\begin{equation}
H_\mathrm{KLRS}(p) \equiv \mathsf{IsCrossoverEW}(p),
\end{equation}
the assertion that the full thermal-resummed SM is a crossover. The
hypothesis is encoded as a tracked Prop because the lattice result is
not derived in our framework; the strict-LO
\texttt{smBenchmarkParams} of Phase 5z.3 is first-order
(\texttt{sm\_benchmark\_is\_first\_order}, $E = 0.01 > 0$). The full
thermal correction drives $E \to 0$ for the resummed SM at the
physical $m_H$.

We then prove:
\begin{itemize}
\item \texttt{sm\_no\_nu\_R\_ewbg\_blocked}: SM-no-$\nu_R$ fails
\EWBGViable{} unconditionally, regardless of transition order
(consequence of \texttt{ewbg\_forbidden\_if\_wall\_intact} and
\texttt{sm\_no\_nu\_R\_wall\_intact}).
\item \texttt{sm\_with\_3nu\_R\_ewbg\_iff\_first\_order\_strong}:
SM with $3\nu_R$ has \EWBGViable{} iff the transition is
baryogenesis-viable (the wall conjunct discharges via
\texttt{sm\_with\_3nu\_R\_wall\_cracks}).
\item \texttt{sm\_with\_3nu\_R\_ewbg\_forbidden\_under\_klrs}: under
$H_\mathrm{KLRS}$, even SM with $3\nu_R$ fails \EWBGViable{}.
\end{itemize}

The correctness-push punchline is
\texttt{sm\_no\_nu\_R\_ewbg\_doubly\_forbidden}: under the KLRS
hypothesis, both the SM-no-$\nu_R$ and SM-with-$3\nu_R$ branches fail
\EWBGViable{} via independent failure modes. The two conjuncts are
\emph{not} P2 bundle-redundant: they are statements about \emph{different}
fermion contents (45 vs 48 mod 16) blocked by \emph{different}
mechanisms (wall vs transition). This is the formal statement of
``SM EWBG is doubly forbidden.''

\begin{figure*}
\centering
\includegraphics[width=0.96\textwidth]{figures/fig_ewbg_allowed_region.png}
\caption{EWBG verdict over the (chirality wall $\times$ transition
order) decomposition. \textbf{Left:} 2$\times$2 outcome matrix.
SM-as-is sits at (intact wall + crossover under $H_\mathrm{KLRS}$);
\textbf{doubly forbidden}. SM-with-$3\nu_R$ sits at (cracked wall +
crossover); transition-blocked. BSM target (extra scalars producing
strong first-order EWPT) sits at (cracked wall + first-order strong);
allowed. \textbf{Right:} the $(m_H, E)$ phase diagram with the
KLRS / CFH lattice endpoint at $m_H = 72.4$~GeV. The PDG SM at
$m_H = 125.20$~GeV sits well to the right of the endpoint (overshoot
ratio $\approx 1.73 > 1.5$), placing the resummed transition firmly
in the crossover region. Lean: \texttt{ewbg\_forbidden\_iff\_wall\_
intact\_or\_not\_viable}, \texttt{sm\_no\_nu\_R\_ewbg\_doubly\_
forbidden}, \texttt{sm\_klrs\_overshoot\_ratio\_gt\_threshold}.}
\label{fig:ewbg-allowed-region}
\end{figure*}

%% ================================================================
\section{Formal Verification Summary}
\label{sec:formal-verification}

The Wave 2 formalization comprises one new Lean module:
\begin{center}
\begin{tabular}{ll}
\toprule
Module & Theorems \\
\midrule
\texttt{EWBaryogenesisChiralityWall.lean} & 16 \\
\bottomrule
\end{tabular}
\end{center}
plus three definitions
(\texttt{sphaleronSuppression}, \texttt{SphaleronsDecoupled},
\texttt{EWBGViable}), two predicate definitions (\texttt{WallIntact},
\texttt{WallCracked}), and one tracked-Prop hypothesis
(\texttt{H\_KLRS\_SM\_Crossover}). The module is zero-sorry,
introduces zero new axioms, and adds zero heartbeat overrides
(Phase 3 pipeline invariant 10). Cross-bridges to four prior modules
in proof bodies:
\begin{itemize}
\item Phase 5e \texttt{Z16AnomalyComputation.three\_gen\_anomalous}
+ \texttt{sm\_anomaly\_with\_nu\_R}: chirality-wall integrity at
SM-no-$\nu_R$ and SM-with-$\nu_R$ fermion contents.
\item Phase 5p \texttt{GaugingStep.gauging\_requires\_z16\_cancellation}:
gauging step requires anomaly cancellation (referenced in §5.1).
\item Phase 5a Wave 4 \texttt{ChiralityWallMaster}: three-pillar
synthesis.
\item Phase 5z.3 \texttt{EWPhaseTransition.crossover\_excludes\_
baryogenesis}: crossover blocks sphaleron decoupling.
\end{itemize}

A Python companion \texttt{src/ew\_baryogenesis/} (3 modules, 23-test
pytest suite, $100\%$ pass) wraps the canonical formulas
\texttt{sphaleron\_suppression}, \texttt{chirality\_wall\_blocks\_ewbg},
and \texttt{ewbg\_viable} (in \texttt{src/core/formulas.py}) plus the
structured \texttt{SMEWBGVerdict} dataclass; the Lean and Python
predicate evaluations agree across all SM benchmarks. EWBG-related
constants live in \texttt{EWBG\_PARAMS} in \texttt{src/core/constants.py}
(sphaleron decoupling threshold, KLRS endpoint, SM Higgs mass,
$\Zsixteen$ anomaly representatives).

%% ================================================================
\section{Outlook}
\label{sec:outlook}

EW baryogenesis is doubly forbidden in the SM: by the chirality-wall
obstruction without right-handed neutrinos, and by the crossover
nature of the resummed-thermal SM electroweak transition with or
without right-handed neutrinos. The dispatch is to leptogenesis or
BSM electroweak Higgs sectors with extra scalars producing a strong
first-order transition.

For leptogenesis, the natural connection is the Phase 5z Wave 2
\texttt{MajoranaRung.lean} sterile-neutrino seesaw: with
$M_R \sim 10^{14}$~GeV from the substrate-anchored hypothesis
\texttt{H\_MR\_FromADWSubstrate\_BCS\_LNV} (constrained by Antusch
\textit{et al.}~\cite{Antusch2003} lepton-number-symmetry obstruction),
out-of-equilibrium decays of right-handed neutrinos
generate a lepton asymmetry that is converted to baryon asymmetry by
electroweak sphalerons \emph{above} the EW phase transition (where
sphalerons are unsuppressed). The crossover transition is then no
longer a problem; the wall is cracked by the $\nu_R$ completion;
leptogenesis succeeds at high $T$.

For BSM EWBG, the requirement is the cracked-wall + first-order-strong
quadrant of Fig.~\ref{fig:ewbg-allowed-region}. Two-Higgs-doublet
models, MSSM (with light stops), or extended-scalar singlets can
produce a strong first-order transition by raising the cubic
coefficient at the physical $m_H$~\cite{ProfumoRamsey2007}. The
Lean formalization of this paper is robust to the specific BSM
mechanism: the biconditional \eqref{eq:biconditional} continues to
hold, the wall must crack (anomaly cancellation extends naturally to
BSM with vector-like additions), and the transition order is the
load-bearing input. Phase 6e (nonlinear stress-energy) inherits the
verdict and may need to choose between leptogenesis and BSM EWBG as
the program's preferred baryogenesis mechanism.

%% ================================================================
\begin{thebibliography}{99}

\bibitem{WangWanYou2022}
J.~Wang, X.-G.~Wen, and S.-T.~Yau,
\emph{Anomaly and cobordism constraints beyond the Standard Model:
Topological force},
Phys.~Rev.~D~\textbf{106}, 045013 (2022); arXiv:2112.14765.

\bibitem{Wang2020}
J.~Wang,
\emph{Anomaly and cobordism constraints beyond the Standard Model},
Phys.~Rev.~Research~\textbf{2}, 013189 (2020); arXiv:1910.14664;
DOI:~\href{https://doi.org/10.1103/PhysRevResearch.2.013189}{10.1103/PhysRevResearch.2.013189}.

\bibitem{KLRS1996}
K.~Kajantie, M.~Laine, K.~Rummukainen, and M.~Shaposhnikov,
\emph{Is there a hot electroweak phase transition at $m_H \gtrsim m_W$?},
Phys.~Rev.~Lett.~\textbf{77}, 2887 (1996); arXiv:hep-ph/9605288;
DOI:~\href{https://doi.org/10.1103/PhysRevLett.77.2887}{10.1103/PhysRevLett.77.2887}.

\bibitem{CsikorFodorHeitger1999}
F.~Csikor, Z.~Fodor, and J.~Heitger,
\emph{Endpoint of the hot electroweak phase transition},
Phys.~Rev.~Lett.~\textbf{82}, 21 (1999); arXiv:hep-ph/9809291;
DOI:~\href{https://doi.org/10.1103/PhysRevLett.82.21}{10.1103/PhysRevLett.82.21}.

\bibitem{CohenKaplanNelson1993}
A.~G.~Cohen, D.~B.~Kaplan, and A.~E.~Nelson,
\emph{Progress in electroweak baryogenesis},
Annu.~Rev.~Nucl.~Part.~Sci.~\textbf{43}, 27 (1993);
DOI:~\href{https://doi.org/10.1146/annurev.ns.43.120193.000331}{10.1146/annurev.ns.43.120193.000331}.

\bibitem{Gavela1994}
M.~B.~Gavela, P.~Hernandez, J.~Orloff, and O.~Pene,
\emph{Standard Model CP-violation and baryon asymmetry},
Mod.~Phys.~Lett.~A~\textbf{9}, 795 (1994); arXiv:hep-ph/9312215.

\bibitem{Huet1994}
P.~Huet and E.~Sather,
\emph{Electroweak baryogenesis and Standard Model CP violation},
Phys.~Rev.~D~\textbf{51}, 379 (1995); arXiv:hep-ph/9404302.

\bibitem{ProfumoRamsey2007}
S.~Profumo, M.~J.~Ramsey-Musolf, and G.~Shaughnessy,
\emph{Singlet Higgs phenomenology and the electroweak phase transition},
JHEP \textbf{08} (2007) 010; arXiv:0705.2425.

\bibitem{Antusch2003}
S.~Antusch, J.~Kersten, M.~Lindner, and M.~Ratz,
\emph{Running neutrino masses, mixings and CP phases: Analytical
results and phenomenological consequences},
Nucl.~Phys.~B~\textbf{674}, 401 (2003); arXiv:hep-ph/0305273.

\bibitem{PDG2024}
S.~Navas \textit{et al.}~(Particle Data Group),
\emph{Review of particle physics},
Phys.~Rev.~D~\textbf{110}, 030001 (2024).

\end{thebibliography}

\end{document}
